\documentclass[11pt]{article}
\usepackage[T1]{fontenc}
\usepackage[utf8]{inputenc}
\usepackage{lmodern}
\usepackage[margin=1.05in]{geometry}
\usepackage{amsmath,amssymb,amsthm,mathtools}
\usepackage{microtype}
\usepackage[hidelinks]{hyperref}
\usepackage{enumitem}
\usepackage{booktabs}
\usepackage{tabularx}
\usepackage{xcolor}

\newtheorem{theorem}{Theorem}[section]
\newtheorem{proposition}[theorem]{Proposition}
\newtheorem{lemma}[theorem]{Lemma}
\newtheorem{corollary}[theorem]{Corollary}
\theoremstyle{definition}
\newtheorem{definition}[theorem]{Definition}
\newtheorem{remark}[theorem]{Remark}

\DeclareMathOperator{\lcm}{lcm}
\newcommand{\N}{\mathbb N}
\newcommand{\Q}{\mathbb Q}
\newcommand{\Z}{\mathbb Z}
\newcommand{\Pk}{\mathcal P}
\newcommand{\Bset}{\mathcal B}
\newcommand{\Ectrl}{E_{\mathrm{ctrl}}}
\newcommand{\Egad}{E_{\mathrm{res}}}
\newcommand{\Emass}{E_{\mathrm{mass}}}
\newcommand{\Qctrl}{Q_{\mathrm{ctrl}}}
\newcommand{\QE}{Q_E}
\newcommand{\sctrl}{\sigma_{\mathrm{ctrl}}}
\newcommand{\sE}{\sigma_E}
\newcommand{\normdist}[1]{\left\lVert #1\right\rVert_{\mathbb R/\mathbb Z}}
\newcommand{\status}[1]{\textsf{#1}}

\title{Distinct Semiprime Egyptian Fractions for Rationals with Squarefree Denominator}
\author{Yuren Tang}
\date{Clean recovery manuscript, 18 July 2026}

\begin{document}
\maketitle

\begin{abstract}
We prove that a positive rational number is a finite sum of reciprocals of
\emph{distinct squarefree semiprimes} if and only if its reduced denominator is
squarefree.  The only load-bearing external analytic theorem in the paper proof
is the prime number theorem.  A self-contained Abel-summation argument gives
\[
  \sum_{2^k\le p<2^{k+1}}\frac1p
  =\frac1k+o\!\left(\frac1k\right),
\]
and this local law yields both eventual prime-cardinality supply in each dyadic
block and a fixed reciprocal-prime mass on an inclusive long block window.

The remaining argument is a finite weighted circle-method construction.  A
control graph encodes residue assignments, a mass pool supplies the required
Bernoulli load, and denominator-sensitive reservoirs damp the frequencies that
remain invisible to the block residues.  We separate the global control
partition, the block-minor fibre estimate, squarefree-CRT reservoir damping,
and the final constant-first/scale-last parameter certificate.

The paper proof and the current formal proof have different analytic backends.
The immutable Lean release \texttt{v0.0.3}, archived under DOI
\texttt{10.5281/zenodo.20767390}, remains machine checked relative to two exact
Rosser--Schoenfeld inputs.  The PNT-to-local-law bridge in this article is not
claimed to be part of that released Lean theorem.
\end{abstract}

\tableofcontents

\section{Statement and proof architecture}

An integer is a \emph{squarefree semiprime} if it is a product $pq$ of two
distinct primes.  Equivalently, its arithmetic functions satisfy
$\omega(n)=\Omega(n)=2$.

\begin{theorem}[Erd\H{o}s Problem 306]\label{thm:headline}
Let $q>0$ be rational and write $q=a/b$ in lowest terms.  Then $q$ is a finite
sum of reciprocals of distinct squarefree semiprimes if and only if $b$ is
squarefree.
\end{theorem}

The necessity is immediate but fixes the exact scope.  If $S$ is a finite set
of squarefree integers, then the reduced denominator of
$\sum_{n\in S}1/n$ divides $\lcm(S)$, which is squarefree.

The sufficiency direction has four layers.
\begin{enumerate}[label=\textbf{Layer \arabic*.}]
\item PNT is converted by partial summation into a reciprocal-prime law on one
      dyadic block.
\item The local law gives the two eventual supplies consumed by the
      construction.
\item A finite spectral construction gives an avoiding representation of
      $1/b$ for every squarefree $b\ge3$.
\item Small denominators and disjoint iteration give arbitrary numerators.
\end{enumerate}
This is a mathematical principle/handoff organization, not a traversal of a
particular source tree.

\section{Avoiding representations and arithmetic closure}

For a finite set $T\subset\N$, write $\operatorname{Rep}_T(r)$ if there is a
finite set $S\subset\N$ such that every $n\in S$ is a squarefree semiprime,
$S\cap T=\varnothing$, and $\sum_{n\in S}1/n=r$.

\begin{lemma}[Numerator reduction]\label{lem:numerator-reduction}
If $\operatorname{Rep}_T(1/b)$ holds for every finite $T$, then
$\operatorname{Rep}_T(a/b)$ holds for every $a\ge1$ and every finite $T$.
\end{lemma}
\begin{proof}
Induct on $a$.  At each step apply the unit-numerator assertion with obstruction
set enlarged by all denominators already chosen.  The resulting finite sets are
disjoint, so their reciprocal sums add.
\end{proof}

\begin{lemma}[Small denominators]\label{lem:small-b}
If the avoiding unit-numerator assertion holds for every squarefree $b\ge3$,
then it holds for all positive squarefree $b$.
\end{lemma}
\begin{proof}
For $b=2$, choose disjoint avoiding representations of $1/3$ and $1/6$ and use
$1/2=1/3+1/6$.  For $b=1$, use disjoint representations of $1/2$, $1/3$, and
$1/6$ together with $1=1/2+1/3+1/6$.
\end{proof}

\section{PNT and the local reciprocal-prime law}\label{sec:PNT}

Let $\pi(x)$ be the number of primes at most $x$.  We use the prime number
theorem in the standard form
\begin{equation}\label{eq:PNT}
  \pi(x)\sim\frac{x}{\log x}\qquad(x\to\infty).
\end{equation}
A directly checkable modern statement is given by Soundararajan
\cite{Soundararajan2006}; Ingham, Theorem~23, is a standard monograph locator
\cite{Ingham1990}.  Hadamard's 1896 paper is cited for historical provenance
\cite{Hadamard1896}.

\begin{proposition}[Local multiplicative-block law]\label{prop:local-law}
The prime number theorem implies
\begin{equation}\label{eq:local-law-x}
  \sum_{x\le p<2x}\frac1p
  =\frac{\log2}{\log x}+o\!\left(\frac1{\log x}\right).
\end{equation}
Consequently, for
\[
 \Pk_k=\{p\text{ prime}:2^k\le p<2^{k+1}\},\qquad
 A_k=\sum_{p\in\Pk_k}\frac1p,
\]
one has
\begin{equation}\label{eq:local-law-k}
  A_k=\frac1k+o\!\left(\frac1k\right).
\end{equation}
\end{proposition}
\begin{proof}
First use the interval $(x,2x]$.  Abel summation for the prime-counting step
function and $f(t)=1/t$ gives
\begin{equation}\label{eq:abel}
 \sum_{x<p\le2x}\frac1p
 =\frac{\pi(2x)}{2x}-\frac{\pi(x)}x
  +\int_x^{2x}\frac{\pi(t)}{t^2}\,dt.
\end{equation}
Put
\[
 \varepsilon(x)=\sup_{t\ge x}
 \left|\frac{\pi(t)\log t}{t}-1\right|.
\]
By PNT, $\varepsilon(x)\to0$, and uniformly for $t\in[x,2x]$,
$\pi(t)=t(1+O(\varepsilon(x)))/\log t$.  Hence
\begin{align*}
 \int_x^{2x}\frac{\pi(t)}{t^2}\,dt
 &=\int_x^{2x}\frac{dt}{t\log t}
   +O\!\left(\varepsilon(x)\int_x^{2x}\frac{dt}{t\log t}\right)\\
 &=\log\frac{\log(2x)}{\log x}
   +o\!\left(\frac1{\log x}\right)\\
 &=\frac{\log2}{\log x}
   +O\!\left(\frac1{\log^2x}\right)
   +o\!\left(\frac1{\log x}\right).
\end{align*}
The boundary term in \eqref{eq:abel} is
\[
 \frac1{\log(2x)}-\frac1{\log x}
 +o\!\left(\frac1{\log x}\right)
 =o\!\left(\frac1{\log x}\right).
\]
Changing $(x,2x]$ to $[x,2x)$ changes the sum by
$O(1/x)=o(1/\log x)$.  At $x=2^k$, $k\ge2$, both endpoints are composite, so
the conventions agree exactly.  Since $\log(2^k)=k\log2$, the dyadic formula
follows.
\end{proof}

\begin{remark}[Why ordinary Mertens is not the local premise]
The relation $\sum_{p\le x}1/p=\log\log x+B+o(1)$ is sufficient for a fixed
power window, but subtraction on one dyadic block leaves an uncontrolled
$o(1)$ error while the desired main term is only of order $1/\log x$.  Thus it
is background, not a second load-bearing theorem here; see Mertens
\cite{Mertens1874}.
\end{remark}

\section{Eventual structural supply}\label{sec:supply}

The construction chooses its bottom scale after every fixed constant and finite
obstruction.  It therefore consumes eventual, rather than fixed-cutoff,
interfaces.

\begin{definition}[Eventual dyadic density $\mathbf D_{\mathrm{ev}}$]\label{def:D}
There exists $K_D$ such that for every $k\ge K_D$,
\begin{equation}\label{eq:D}
 |\Pk_k|\ge\frac{2^k}{2\log(2^k)}.
\end{equation}
\end{definition}

\begin{definition}[Eventual reciprocal window $\mathbf M_{\mathrm{ev}}$]\label{def:M}
There exists $K_M$ such that for every $k_0\ge K_M$,
\begin{equation}\label{eq:M}
 \sum_{k=k_0}^{3k_0}A_k\ge\frac{21}{20}.
\end{equation}
The upper block index is inclusive, so the associated prime interval is
$[2^{k_0},2^{3k_0+1})$.
\end{definition}

\begin{proposition}[The local law supplies both interfaces]\label{prop:local-to-supply}
Equation \eqref{eq:local-law-k} implies
$\mathbf D_{\mathrm{ev}}$ and $\mathbf M_{\mathrm{ev}}$.
\end{proposition}
\begin{proof}
For $p\in\Pk_k$, $1/p\le2^{-k}$, whence
\[
 A_k\le\frac{|\Pk_k|}{2^k},\qquad |\Pk_k|\ge2^kA_k.
\]
Eventually $A_k\ge0.9/k$, and $0.9>1/(2\log2)$, which gives
\eqref{eq:D}.  Also eventually $A_k\ge0.98/k$, so
\[
 \sum_{k=k_0}^{3k_0}A_k
 \ge0.98\sum_{k=k_0}^{3k_0}\frac1k.
\]
The harmonic window tends to $\log3$, and
$0.98\log3>21/20$, proving \eqref{eq:M}.
\end{proof}

The two supplies have different roles.  Density populates individual control
blocks and high-prime reservoirs; reciprocal mass supplies a pool of semiprime
reciprocals.  Neither statement alone substitutes for the other.

\begin{theorem}[Structural avoiding construction]\label{thm:structural-construction}
Assume $\mathbf D_{\mathrm{ev}}$ and $\mathbf M_{\mathrm{ev}}$.  If $b\ge3$ is
squarefree and $T\subset\N$ is finite, then $\operatorname{Rep}_T(1/b)$ holds.
\end{theorem}

\begin{corollary}\label{cor:structural-E306}
The prime number theorem implies the sufficiency direction of
Theorem~\ref{thm:headline}.
\end{corollary}
\begin{proof}
Use Propositions~\ref{prop:local-law} and \ref{prop:local-to-supply}, then
Theorem~\ref{thm:structural-construction}, Lemma~\ref{lem:small-b}, and
Lemma~\ref{lem:numerator-reduction}.
\end{proof}

The remainder of the paper proves Theorem~\ref{thm:structural-construction}.
Its constants are chosen first and $k_0$ last, so replacing fixed explicit
cutoffs by eventual thresholds only enlarges the final lower bound for $k_0$.

\section{Finite spectral selection}\label{sec:spectral}

Let $E$ be a finite set of positive integers, let $L$ be a common multiple of
$E$, let $q_0\in\Z/L\Z$, and let $\theta_e\in[0,1]$.  Give a subset
$A\subseteq E$ the product weight
\[
 w(A)=\prod_{e\in A}\theta_e\prod_{e\notin A}(1-\theta_e).
\]
The weighted number of subsets satisfying
$\sum_{e\in A}L/e\equiv q_0\pmod L$ is
\[
 W_{L,q_0}(E,\theta)=
 \sum_{\substack{A\subseteq E\\
 \sum_{e\in A}L/e\equiv q_0\ (\mathrm{mod}\ L)}}w(A).
\]
Finite Fourier orthogonality gives
\begin{equation}\label{eq:finite-fourier}
 L W_{L,q_0}(E,\theta)
 =\sum_{h=0}^{L-1}
   \prod_{e\in E}\bigl((1-\theta_e)+\theta_e e^{2\pi i h/e}\bigr)
   e^{-2\pi i hq_0/L}.
\end{equation}

\begin{lemma}[Finite spectral selection]\label{lem:spectral-selection}
Partition the frequencies as $S_M\dot\cup S_m$.  If the main sum is real and
\[
 \Re\sum_{h\in S_M}F(h)>\sum_{h\in S_m}|F(h)|,
\]
where $F(h)$ is the summand in \eqref{eq:finite-fourier}, then
$W_{L,q_0}(E,\theta)>0$.
\end{lemma}
\begin{proof}
The real part of the complete Fourier sum is positive.  Since all subset
weights are nonnegative, at least one admissible subset has positive weight.
\end{proof}

For the reciprocal problem take $q_0=L/b$.  If every $e$ divides $L$ and
\begin{equation}\label{eq:no-wrap}
 \sum_{e\in E}\frac1e<1,
\end{equation}
then the congruence is equivalent to
$\sum_{e\in A}1/e=1/b$, because both sides lie in $[0,1)$.

\section{Dyadic control geometry}\label{sec:control}

Fix squarefree $b\ge3$ and finite $T$.  Choose a large bottom scale $k_0$, a
terminal scale $K$, and prime sets
$P_k\subseteq\Pk_k$ for $k_0\le k\le K$.  Put
$\Bset=\bigcup_{k=k_0}^K P_k$ and choose every prime in $\Bset$ larger than
$b$ and every element of $T$.

A control graph on $\Bset$ contains sufficiently many pairs inside each block
and selected pairs between consecutive blocks.  Let
\[
 \Ectrl=\{pq:\{p,q\}\text{ is a control pair}\}.
\]
Unique factorization makes these denominators distinct.  For an assignment
$a=(a_p)_{p\in\Bset}\in\prod_{p\in\Bset}\Z/p\Z$, let $x_{pq}(a)$ be the
centered CRT representative determined by $(a_p,a_q)$ and set
\begin{equation}\label{eq:Qctrl}
 \Qctrl(a)=\sum_{\{p,q\}\in\mathcal C}
 \left(\frac{x_{pq}(a)}{pq}\right)^2,
 \qquad
 \sctrl^2=\sum_{\{p,q\}\in\mathcal C}\frac1{(pq)^2}.
\end{equation}
For a frequency $h$, the exact CRT identity gives
\[
 \normdist{h/(pq)}=\frac{|x_{pq}(a(h))|}{pq},
\]
and therefore
\begin{equation}\label{eq:control-full-energy}
 \Qctrl(a(h))\le \QE(h):=\sum_{e\in E}\normdist{h/e}^2.
\end{equation}

The graph is chosen so that, once $k_0$ is sufficiently large,
\begin{align}
 \sctrl&\ge\frac1{c_\sigma k_0 2^{k_0}},\label{eq:sctrl-lower}\\
 \sum_{e\in\Ectrl}\frac1e&<\frac{3}{4b}.\label{eq:control-load}
\end{align}
The first estimate comes from the numerous bottom-block pairs and the second
from quadratic growth of semiprime denominators.  Neither uses reciprocal mass.

For $C\ge1$ define the main assignment set
\begin{equation}\label{eq:assignment-main}
 \mathfrak M(C)=\left\{a:\exists m\in\Z,\ |m|\le C/\sctrl,
 \ a_p\equiv m\pmod p\ \forall p\in\Bset\right\}.
\end{equation}

\section{The global control partition}\label{sec:global}

\begin{theorem}[Global control partition]\label{thm:global-partition}
For every $c>0$ and $\eta>0$ there exist $k_*(c,\eta)$ and
$C_{\mathrm{tail}}>0$ such that, for every admissible block system with
$k_0\ge k_*$ and every $C\ge1$,
\begin{equation}\label{eq:global-partition}
 \sum_{a\notin\mathfrak M(C)}e^{-c\Qctrl(a)}
 \le\frac{\eta+C_{\mathrm{tail}}e^{-cC^2/2}}{\sctrl}.
\end{equation}
\end{theorem}

We record the four handoffs in the proof.  They are also the natural units for
independent review.

\subsection{Cold blocks and finite encoding}
For a block $P_k$, write $Q_k(a)$ for its local control energy.  The control
pairs provide a forcing threshold $R_k$ and the pairs between consecutive
blocks provide a boundary threshold $\Pi_k$.

\begin{lemma}[Cold-block structure]\label{lem:cold-block}
For large $k_0$, if $Q_k(a)<R_k$, then one integer label $\ell_k$ agrees with
$a_p$ modulo all but an energy-controlled exceptional set of primes in $P_k$.
The number of labels and exceptional sets with energy at most $R$ is bounded by
an exponential charge $e^{\varepsilon R}$ times a polynomial label window.
\end{lemma}
\begin{proof}[Mechanism]
Low phase energy on many pairs forces their centered CRT representatives into a
short interval.  Uniqueness in that interval decodes one integer fingerprint.
If two substantial residue classes survived, their cross pairs would contribute
a quadratic amount of energy.  Each exceptional prime therefore pays a fixed
energy charge, and summing the exceptional choices uses
$\sum_A\prod_{j\in A}u_j=\prod_j(1+u_j)$.
\end{proof}

\begin{lemma}[Adjacent-label mismatch]\label{lem:boundary-penalty}
Distinct dominant labels on consecutive cold blocks force energy at least
$\Pi_k$ at their common boundary.
\end{lemma}
\begin{proof}[Mechanism]
Cross-block control pairs determine centered CRT representatives from the two
labels.  If the labels differ, sufficiently many representatives are nonzero
and of controlled size.  Squaring and summing gives the boundary floor; removed
exceptional primes have already paid their energy charge.
\end{proof}

An assignment of energy at most $R$ is encoded by its hot blocks, its cold
boundaries with a label change, dyadic energy shells, one label on each cold
segment, and the residual exceptional fibre.  Every hot block and mismatch
boundary spends a definite part of the energy budget.

\begin{lemma}[Level-set estimate]\label{lem:level-set}
For every $0<\varepsilon<1$ there are $A>0$ and $k_\varepsilon$ such that, for
$k_0\ge k_\varepsilon$ and $R\ge1$,
\begin{equation}\label{eq:level-set}
 \#\{a:\Qctrl(a)\le R\}
 \le e^{AJ}e^{8\varepsilon R}
 \left(1+\frac{\sqrt R}{\sctrl}\right),
\end{equation}
where $J$ is the number of blocks.
\end{lemma}
\begin{proof}
Sum the encoding fibres over hot blocks, mismatch boundaries, shells, and cold
segments.  The weighted subset sums for hot blocks and boundaries contribute
at most $e^{2\varepsilon R+J}$ each; shell and exception data contribute
$e^{4\varepsilon R+A_0J}$.  Once the first cold-segment label is fixed, labels
propagate except across charged boundaries.  The first label has size
$O(\sqrt R/\sctrl)$ because a diagonal label $m$ already contributes
$m^2\sctrl^2$.
\end{proof}

\subsection{Localization and Laplace absorption}
Let $\mathcal F_0=\min(R_{k_0},\Pi_{k_0})$ and define
\[
 \mathfrak D(C)=\{a:\exists m\in\Z,\ a_p\equiv m\pmod p\ \forall p,
 |m|>C/\sctrl,\ \Qctrl(a)=m^2\sctrl^2\}.
\]

\begin{lemma}[Localization dichotomy]\label{lem:localization}
For sufficiently large $k_0$,
\[
 a\notin\mathfrak M(C)\Longrightarrow
 \Qctrl(a)\ge\mathcal F_0\quad\text{or}\quad a\in\mathfrak D(C).
\]
\end{lemma}
\begin{proof}
A hot block or a changed cold label gives the first alternative.  Otherwise all
blocks are cold, the low-energy merger removes the exceptions, and all labels
coincide with one integer $m$.  If $|m|\le C/\sctrl$ the assignment is main;
otherwise it lies in $\mathfrak D(C)$.
\end{proof}

\begin{lemma}[Mass above the forcing floor]\label{lem:floor-laplace}
For every $c>0$ and $\eta>0$, after increasing $k_0$,
\begin{equation}\label{eq:floor-laplace}
 \sum_{\Qctrl(a)\ge\mathcal F_0}e^{-c\Qctrl(a)}
 \le\frac\eta{\sctrl}.
\end{equation}
\end{lemma}
\begin{proof}
Choose $\varepsilon$ with $8\varepsilon<c$ and an intermediate exponent
$\bar c\in(8\varepsilon,c)$.  Apply Lemma~\ref{lem:level-set} on integer
energy shells.  The resulting full Laplace sum is at most
$e^{AJ}K_0(1+1/\sctrl)$.  On the sector above $\mathcal F_0$, extract
$e^{-(c-\bar c)\mathcal F_0}$.  The forcing floor grows sufficiently quickly
with $k_0$ to absorb the fixed entropy factor and leave coefficient $\eta$.
\end{proof}

\begin{lemma}[Diagonal Gaussian tail]\label{lem:diagonal-tail}
For every $c>0$ there is $C_{\mathrm{tail}}>0$ such that
\begin{equation}\label{eq:diagonal-tail}
 \sum_{a\in\mathfrak D(C)}e^{-c\Qctrl(a)}
 \le\frac{C_{\mathrm{tail}}e^{-cC^2/2}}{\sctrl}.
\end{equation}
\end{lemma}
\begin{proof}
A diagonal assignment is determined by its integer label, so the sum injects
into
\[
 \sum_{|m|>C/\sctrl}e^{-cm^2\sctrl^2}.
\]
Split the exponent in half: one half supplies $e^{-cC^2/2}$, while the
remaining integer Gaussian sum is $O_c(1/\sctrl)$.
\end{proof}

Theorem~\ref{thm:global-partition} follows by applying the localization
dichotomy and adding the last two estimates.  The factor $e^{AJ}$ has $A$
fixed before the final scale is chosen; this order is essential.

\section{Reciprocal mass and the weighted edge family}\label{sec:mass}

Put
\[
 \Pk(k_0)=\bigcup_{k=k_0}^{3k_0}\Pk_k,\qquad
 S_1=\sum_{p\in\Pk(k_0)}\frac1p,\qquad
 S_2=\sum_{p\in\Pk(k_0)}\frac1{p^2}.
\]
For products of two distinct primes,
\begin{equation}\label{eq:pair-mass-id}
 \sum_{\substack{p,q\in\Pk(k_0)\\p<q}}\frac1{pq}
 =\frac{S_1^2-S_2}{2}.
\end{equation}
Condition $\mathbf M_{\mathrm{ev}}$ gives $S_1\ge21/20$, while
$S_2\le(2^{k_0}-1)^{-1}$.  Hence, after increasing $k_0$,
\begin{equation}\label{eq:pair-mass-half}
 \sum_{p<q}\frac1{pq}\ge\frac12.
\end{equation}

The denominator family is a disjoint union
$E=\Ectrl\,\dot\cup\,\Emass\,\dot\cup\,\Egad$.
The control and reservoir sets are fixed before the mass batch and satisfy
\begin{equation}\label{eq:base-load}
 \sum_{e\in\Ectrl\cup\Egad}\frac1e<\frac{3}{2b}.
\end{equation}
After removing $T$, the control products, and the reservoir products from the
pair pool, the remaining load still exceeds the required residual target at
large scale.

\begin{lemma}[Greedy reciprocal window]\label{lem:greedy-window}
If positive reals $(x_i)$ have total at least $t$ and each $x_i<g$, then some
subfamily has sum in $[t,t+g)$.
\end{lemma}
\begin{proof}
Take a subfamily minimal under inclusion among those with sum at least $t$.
Removing any selected term drops the sum below $t$.
\end{proof}

Apply the lemma to the residual target
$t=3/(2b)-\sum_{e\in\Ectrl\cup\Egad}1/e$ and gap $g=3/(2b)$.  It gives a
mass batch such that
\begin{equation}\label{eq:load-window}
 \frac{3}{2b}\le\sum_{e\in E}\frac1e<\frac3b.
\end{equation}
All denominators are distinct squarefree semiprimes, avoid $T$, and divide
\begin{equation}\label{eq:period}
 L=b\prod_{p\in\Bset}p.
\end{equation}
Let $\Lambda=\sum_{e\in E}1/e$ and use the uniform weight
\begin{equation}\label{eq:theta}
 \theta=\frac{1/b}{\Lambda}.
\end{equation}
Then
\begin{equation}\label{eq:theta-window}
 \frac13<\theta\le\frac23,\qquad
 \sum_{e\in E}\frac\theta e=\frac1b.
\end{equation}
The upper load bound gives \eqref{eq:no-wrap}.  Finally, for an absolute
$S_{\mathrm{load}}$,
\begin{equation}\label{eq:square-load}
 \sum_{e\in E}\frac1{e^2}\le S_{\mathrm{load}}\sctrl^2.
\end{equation}
Thus, with
$\sE^2=\sum_{e\in E}\theta(1-\theta)/e^2$, there is $K_\sigma\ge1$ such that
\begin{equation}\label{eq:variance-comparison}
 \sqrt{\frac29}\,\sctrl\le\sE\le K_\sigma\sctrl.
\end{equation}

\section{Main Fourier arcs}\label{sec:main}

Choose $C\ge3$ and put
\begin{equation}\label{eq:N}
 N=\left\lceil\frac C{\sctrl}\right\rceil.
\end{equation}
After enlarging $k_0$,
\begin{equation}\label{eq:window-geometry}
 2N<2^{2k_0},\qquad 2N+1\le L,\qquad \frac1{\sE}\le N.
\end{equation}
For $|m|\le N$, let
\[
 T_m=\prod_{e\in E}\bigl((1-\theta)+\theta e^{2\pi i m/e}\bigr)
 e^{-2\pi i m/b}.
\]
The exact expected-mass identity in \eqref{eq:theta-window} cancels the linear
term in the logarithm.  A uniform third-order expansion gives
\begin{equation}\label{eq:main-expansion}
 \log T_m=-2\pi^2m^2\sE^2+\delta_m,
 \qquad |\delta_m|\le C_3\sum_{e\in E}|m/e|^3.
\end{equation}
The scale inequality places every $m/e$ in the Taylor disk, and
\eqref{eq:square-load} makes the remainder at most $1/10$.  Hence
\begin{equation}\label{eq:per-main-lower}
 \Re T_m\ge0.8e^{-2\pi^2m^2\sE^2}.
\end{equation}
Conjugate symmetry makes the main sum real.  A finite Gaussian lower bound then
gives
\begin{equation}\label{eq:main-lower}
 \Re\sum_{|m|\le N}T_m\ge\frac{c_3}{\sE},
 \qquad c_3=0.8\frac{e^{-\pi^2/2}}2>0.
\end{equation}

For every frequency,
\[
 |(1-\theta)+\theta e^{2\pi it}|^2
 =1-4\theta(1-\theta)\sin^2(\pi t).
\]
Using $\theta\in[1/3,2/3]$ and the elementary lower bound for
$\sin^2(\pi t)$ in terms of $\normdist t^2$ gives
\begin{equation}\label{eq:fourier-energy}
 |F(h)|\le e^{-(16/9)\QE(h)}.
\end{equation}

\section{The block-minor fibre estimate}\label{sec:block-minor}

Map $h\pmod L$ to $a(h)=(h\pmod p)_{p\in\Bset}$.  The block-minor sector is
$a(h)\notin\mathfrak M(C)$.  Since
$L=b\prod_{p\in\Bset}p$, fixing the block residues leaves at most $b$
frequencies modulo $L$.  Write
$Q_{\mathrm{extra}}(h)=\QE(h)-\Qctrl(a(h))\ge0$.  Each fibre satisfies
\begin{equation}\label{eq:fiber-tail}
 \sum_{\substack{h\text{ in block sector}\\a(h)=a}}
 e^{-(16/9)Q_{\mathrm{extra}}(h)}\le b.
\end{equation}
Reindexing by fibres and applying Theorem~\ref{thm:global-partition},
\begin{equation}\label{eq:block-budget}
 \sum_{h\text{ in block sector}}e^{-(16/9)\QE(h)}
 \le\frac{b\bigl(\eta+C_{\mathrm{tail}}e^{-(16/9)C^2/2}\bigr)}{\sctrl}.
\end{equation}
The factor $b$ is the exact CRT fibre multiplicity and is not hidden in
asymptotic notation.

\section{Squarefree CRT and reservoir damping}\label{sec:extra}

The remaining minor frequencies lie in a main assignment fibre and therefore
carry a label $m$ with $|m|\le C/\sctrl\le N$ and
\begin{equation}\label{eq:block-agreement}
 h\equiv m\pmod p\qquad(p\in\Bset).
\end{equation}
A genuine main frequency is the residue of $m$ modulo $L$.  There are at most
$b-1$ other siblings for each label, so
\begin{equation}\label{eq:extra-count}
 |S_{\mathrm{extra}}|\le b(2N+1).
\end{equation}

\begin{lemma}[Squarefree sibling mismatch]\label{lem:sibling-mismatch}
If \eqref{eq:block-agreement} holds but $h\not\equiv m\pmod L$, then some prime
$r\mid b$ satisfies $h\not\equiv m\pmod r$.
\end{lemma}
\begin{proof}
Agreement modulo every prime divisor of squarefree $b$ would imply agreement
modulo $b$.  Together with the block congruences and coprimality, CRT would then
give agreement modulo $L$.
\end{proof}

For every prime $r\mid b$, choose a common set $S$ of $G$ high block primes and
include
\begin{equation}\label{eq:reservoir}
 \Egad=\{rs:r\mid b\text{ prime},\ s\in S\}.
\end{equation}
The geometry gives $2|m|<s$.  For the mismatch prime and every $s\in S$,
\begin{equation}\label{eq:one-gadget}
 \left|(1-\theta)+\theta e^{2\pi ih/(rs)}\right|
 \le\sqrt{1-\frac8{9r^2}}
 \le\sqrt{1-\frac8{9b^2}}=: \beta_b<1.
\end{equation}
Thus $|F(h)|\le\beta_b^G$.  Choose $D_{\mathrm{mp}}>0$ and then $G$ so that
$\beta_b^G\le D_{\mathrm{mp}}$.  Equation \eqref{eq:extra-count} yields
\begin{equation}\label{eq:extra-budget}
 \sum_{h\in S_{\mathrm{extra}}}|F(h)|
 \le b(2N+1)D_{\mathrm{mp}}.
\end{equation}
Condition $\mathbf D_{\mathrm{ev}}$ supplies the $G$ high primes once $k_0$
exceeds its threshold.

\section{Terminal parameter compatibility}\label{sec:terminal}

Let $c_3$ be as in \eqref{eq:main-lower}.  Choose $K_\sigma\ge1$ so that the
upper bound in \eqref{eq:variance-comparison} holds.  The constants are fixed
in the following order.

\begin{enumerate}[label=\textbf{Step \arabic*.}]
\item Put
\begin{equation}\label{eq:eta-choice}
 \eta=\frac{c_3}{4K_\sigma b}.
\end{equation}
The global partition returns $C_{\mathrm{tail}}$ and a scale threshold.
\item Choose $C\ge3$ so large that
\begin{equation}\label{eq:gaussian-quarter}
 bC_{\mathrm{tail}}e^{-(16/9)C^2/2}<\frac{c_3}{4K_\sigma}.
\end{equation}
\item Put
\begin{equation}\label{eq:Dmp}
 D_{\mathrm{mp}}=\frac{c_3}{4K_\sigma b(2C+3)},
\end{equation}
and choose $G$ with $\beta_b^G\le D_{\mathrm{mp}}$.
\item Choose $k_0$ only now, above the finite maximum of all thresholds for the
global partition, $\mathbf D_{\mathrm{ev}}$, $\mathbf M_{\mathrm{ev}}$, the
$G$ high primes, the small base load, avoidance of $T$, the variance estimates,
window geometry, Taylor error, and inverse-square load.
\end{enumerate}

The single terminal budget is
\begin{equation}\label{eq:terminal-budget}
 b\eta+bC_{\mathrm{tail}}e^{-(16/9)C^2/2}
 +bD_{\mathrm{mp}}(2C+3)<\frac{c_3}{K_\sigma}.
\end{equation}
Moreover, since $N=\lceil C/\sctrl\rceil$,
$(2N+1)\sctrl\le2C+3$.  Combining \eqref{eq:block-budget},
\eqref{eq:extra-budget}, and $\sE\le K_\sigma\sctrl$ gives
\begin{equation}\label{eq:minor-beat}
 \sum_{h\in S_m}|F(h)|<\frac{c_3}{\sE}.
\end{equation}
The main sum is real and at least $c_3/\sE$ by \eqref{eq:main-lower}.
Lemma~\ref{lem:spectral-selection} therefore selects $A\subseteq E$ with
$\sum_{e\in A}1/e=1/b$.  Its elements are distinct squarefree semiprimes and
avoid $T$, proving Theorem~\ref{thm:structural-construction}.

\begin{remark}
The order prevents four circularities: $C$ is chosen only after the tail
constant is known; $G$ only after the extra-frequency budget is fixed; $k_0$
does not define a constant it must later absorb; and the mass batch is selected
only after control and reservoir loads are budgeted.
\end{remark}

\section{Formal statement and exact trust boundary}\label{sec:trust}

The released Lean theorem has the form
\begin{align*}
 \forall q\in\Q,\quad &0<q\to\operatorname{Squarefree}(q.\mathrm{den})\to\\
 &\exists k\in\N\ \exists n:\operatorname{Fin}(k+1)\to\N,\quad
 n(0)=1,\quad n\text{ strictly increasing},\\
 &\omega(n_i)=\Omega(n_i)=2\ (1\le i\le k),\qquad
 q=\sum_{i=1}^k\frac1{n_i}.
\end{align*}
The value $n(0)=1$ is a dummy anchor; strict monotonicity makes the remaining
denominators distinct.

\subsection{Paper proof}
The sole load-bearing external analytic theorem in this article is PNT
\eqref{eq:PNT}.  Proposition~\ref{prop:local-law} is proved in the text, and its
eventual consequences feed the construction.  Ordinary reciprocal-prime
Mertens and the explicit Rosser--Schoenfeld inequalities are not premises of
the paper proof.

\subsection{Immutable released evidence}
The current archived formal evidence is
\[
 \texttt{v0.0.3@4582185de1e0e27416e9362e0cc7943c3d2fb4fe},
\]
cited as \cite{TangErdos306v003} and archived under DOI
\texttt{10.5281/zenodo.20767390}.  Its audit reports exactly
\begin{gather*}
 \texttt{propext},\qquad \texttt{Classical.choice},\qquad
 \texttt{Quot.sound},\\
 \texttt{RosserSchoenfeld.rosser\_schoenfeld\_cor3},\\
 \texttt{RosserSchoenfeld.rosser\_schoenfeld\_thm5}.
\end{gather*}
The first three are standard Lean foundations.  The last two are explicit
source-transcribed inputs from Rosser--Schoenfeld \cite{RosserSchoenfeld1962},
Corollary~3, equation~(3.8), p.~69, and Theorem~5,
equations~(3.17)--(3.18), p.~70.  They remain the exact formal backend of the
release.  The PNT argument in this paper does not retroactively alter that
trust boundary.

\subsection{Frozen architecture context}
A later frozen checkpoint,
\[
 \texttt{codex/pushlinter@e55ef359a8b98525f0bac6c7a510fcad94469bff},
\]
separates the public statement, structural analytic inputs, finite spectral
selection, global control, construction certificates, and numerator reduction.
Its mathematical handoff structure informs the exposition above, but it is
unreleased and is not cited as the current formal authority.  Its structural
analytic declarations are not represented here as proved Lean consequences.

\subsection{Review and publication status}
The present text is a paper-level reconstruction, not a new kernel object.
Independent review remains required for the PNT endpoint calculation, the
global-control encoding, the block-minor transfer, the reservoir damping, and
the terminal constant hierarchy.  Completion of this manuscript recovery does
not constitute independent review, publication clearance, or authorization to
submit to arXiv.

\section{Reproducibility, novelty, and disclosure}

A clean checkout of the immutable release can be checked with its pinned Lean
and Mathlib environment using
\begin{verbatim}
cd lean
lake exe cache get
lake build
lake env lean RequestProject/Audit.lean
\end{verbatim}
The final command prints the five-item audit above.  The archived release, not
this manuscript branch or the frozen refactor, is the stable software citation.

The conservative novelty claim is that the archived software contains a
machine-checked proof of the stated Formal Conjectures formulation of Erd\H{o}s
Problem~306, conditional on the two named Rosser--Schoenfeld statements, and
that this manuscript gives a PNT-first mathematical reconstruction of the
proof architecture.  No claim of peer review or independent verification is
made.  The benchmark and historical problem sources are
\cite{FormalConjectures2026,ErdosGraham1980}.

The project used AI-assisted mathematical exploration, Lean development,
source organization, and manuscript drafting.  Responsibility for the
mathematical claims, source selection, and final submitted text remains with
the author.  The final disclosure wording and acknowledgements remain author
decisions.

\appendix
\section{Dependency index for independent review}
The proof may be audited in the following order.
\begin{enumerate}
\item PNT, Abel summation, endpoint conventions, and the local block law.
\item The two eventual supply consequences and the inclusive endpoint.
\item Finite Fourier identity, positivity criterion, and no-wrap conversion.
\item Control CRT identity, variance scale, and reciprocal load.
\item Cold-block decoding, mismatch energy, level-set count, localization,
      high-floor absorption, and the diagonal Gaussian tail.
\item Pair-pool identity, forbidden-load subtraction, greedy window, and
      variance comparison.
\item Main-arc Taylor error and the finite Gaussian lower bound.
\item Exact block-fibre multiplicity $b$ and the reindexing inequality.
\item Squarefree CRT mismatch and product reservoir damping.
\item Constant-first/scale-last parameter compatibility.
\item Small-denominator and numerator reductions.
\item Separate comparison of the paper backend with the released formal audit.
\end{enumerate}

\section{Status vocabulary}
\begin{description}
\item[\status{Paper external}.] A cited theorem used by the article.
\item[\status{Paper proved}.] A derivation supplied in the mathematical text.
\item[\status{Released checked}.] A declaration accepted by the immutable
release's Lean kernel environment.
\item[\status{Source-transcribed formal input}.] An external theorem represented
as a named axiom with an exact source locator.
\item[\status{Frozen architecture context}.] An unreleased checkpoint supplies a
clean name or dependency boundary but not current formal authority.
\item[\status{Review gate}.] Independent mathematical, source, formal, or
publication review remains outstanding.
\end{description}

\bibliographystyle{alpha}
\bibliography{references}
\end{document}
